\documentclass[UTF8,a4paper,11pt,fontset=none]{ctexart}
\IfFileExists{C:/Windows/Fonts/simsun.ttc}{%
  \setCJKmainfont[Path=C:/Windows/Fonts/,FontIndex=0,AutoFakeBold=2]{simsun.ttc}
  \setCJKsansfont[Path=C:/Windows/Fonts/,AutoFakeBold=2]{simhei.ttf}
}{%
  \setCJKmainfont[BoldFont=FandolSong-Bold.otf]{FandolSong-Regular.otf}
  \setCJKsansfont[BoldFont=FandolHei-Bold.otf]{FandolHei-Regular.otf}
}
\usepackage[margin=25mm,top=20mm,bottom=20mm]{geometry}
\usepackage{graphicx,booktabs}
\providecommand{\ReportAssets}{}
\providecommand{\ReportEdition}{研究报告初稿}
\setlength{\parindent}{2em}
\setlength{\parskip}{4pt}
\linespread{1.12}
\pagestyle{plain}
\title{\textbf{UCI HAR 时序活动识别}\\[4pt]\large 四模型基线与 CNN 三种子消融}
\author{SequenceModel-HAR-Lab}
\date{\ReportEdition}
\begin{document}
\maketitle

\begin{abstract}
本报告比较同一 UCI HAR 受试者划分上的 CNN、GRU、Transformer 和 Mamba 四个 Seed42 基线，并通过九组历史运行检验 CNN 中加入通道注意力 SE 的表现。当前记录中 GRU 的测试 Macro-F1 最高。SE 的三种子均值高于原 CNN，但配对变化的方向不一致，均值也略低于等参数量残差 MLP。报告使用已归档指标，结论限于当前实现、种子和划分。
\end{abstract}

\section{数据与评估协议}
使用 UCI HAR 官方预处理的惯性信号，每个窗口包含 128 个时间点和 9 个通道，对应 50 Hz 下的 2.56 秒。最终训练、验证、官方测试分别包含 6,037、1,315、2,947 个窗口，受试者分别为 17、4、9 人，三集合互不重叠。仅用最终训练集拟合逐通道均值和标准差。每个运行依据验证 Macro-F1 选模，同分时优先较低验证损失。表中损失为归档的交叉熵，Macro-F1 是六类 F1 的算术平均。

\section{四模型 Seed42 基线}
CNN 用局部卷积提取窗口特征，GRU 用递归末状态汇总序列，Transformer 用自注意力和时间均值，Mamba 使用官方原始选择性状态空间块和时间均值。表~\ref{tab:baselines} 使用各运行的最佳验证 checkpoint 所对应的历史测试结果。基线名称中的 B/C 与后文消融组别独立。

\begin{table}[ht!]
\centering
\caption{官方测试集基线。轮数为最佳／实际，参数为可训练参数。}
\label{tab:baselines}
\small
\begin{tabular}{lrrrrrl}
\toprule
模型 & Accuracy (\%) & Macro-F1 (\%) & Loss & 参数 & 轮数 & 设备 \\
\midrule
\input{\ReportAssets baseline_rows}
\end{tabular}
\end{table}

前三模型在 CPU 上运行，Mamba 在 Linux RTX 3090 上运行。设备和计时边界不同，耗时不能用来排序四种模型的速度。四个基线各仅有一个种子，无法估计架构表现随初始化或受试者划分变化的幅度。

\noindent\textbf{复现材料}　归档 CSV 为制表和绘图输入，来源哈希见 \texttt{input\_manifest.json}。

\clearpage
\section{CNN 三种子消融}
A 为原始 CNN，B 在时间均值前加入 SE 通道门控，C 在时间均值后加入残差 MLP。A 有 27,270 个参数，B 和 C 各有 31,510 个参数。模块位置和计算方式均有差异，B--C 衡量两个具体增量方案之间的差异。

B42、C42 在训练时保存了共享卷积与分类头的初始权重指纹，以及各自逐轮实际读取的训练批次顺序指纹。A42 沿用历史指标；原训练未保存这两类指纹。后续审计按原始源码和 Seed42 确定性配置重建初始权重，并用原数据加载器回放采样顺序。B42、C42 的直接记录在各自覆盖的轮次内与重建结果一致。A42 的配对依据来自事后重建。

\begin{table}[ht!]
\centering
\caption{九组官方测试结果及 Macro-F1 的三种子汇总。标准差采用样本标准差。}
\label{tab:ablation}
\small
\setlength{\tabcolsep}{4pt}
\begin{tabular}[t]{llrrrr}
\toprule
种子 & 组别 & Accuracy (\%) & Macro-F1 (\%) & Loss & 轮数 \\
\midrule
\input{\ReportAssets ablation_rows}
\end{tabular}
\hspace{1em}
\begin{tabular}[t]{lrr}
\toprule
组别 & 均值 (\%) & 标准差 (pp) \\
\midrule
\input{\ReportAssets summary_rows}
\end{tabular}
\end{table}

\begin{figure}[ht!]
\centering
\includegraphics[width=0.94\linewidth]{\ReportAssets results_figure.pdf}
\caption{Seed42 基线差异与同种子 SE 的配对差值。右图保留全部三个观察值，pp 为百分点。三种子固定同一受试者划分，图中的点反映训练随机性下的观察差异。}
\label{fig:results}
\end{figure}

\section{解释与局限}
图~\ref{fig:results} 中，SE 相对 A 的平均提升主要来自 Seed43；在 Seed42 和 Seed44 上差值均为负。SE 与等参数量残差 MLP 的配对差值也有正有负。三个观察值不足以支持稳定提升或统计显著性，现有结果也无法确定改进来自通道门控机制。进一步检验可预先约定更多种子和受试者划分，并比较活动类别召回与不同受试者上的表现。

\end{document}
